% संपूर्ण मराठी ग्रंथातील प्रकरण 28: विन्यासमीमांसेचे अंकगणितीकरण
% हा संपादनयोग्य स्रोतखंड आहे. संकलनासाठी संपूर्ण स्रोतसंग्रहातील
% अक्षररूपे, पूर्वभाग, चिन्हव्याख्या आणि आकृती-स्रोत आवश्यक आहेत.
% मूळ: Open Logic Project; मजकूर आणि रूपांतर CC BY 4.0.
% यंत्रानुवाद, दुरुस्ती आणि तपासणी: OpenAI Codex —
% GPT-5.6 Sol आणि GPT-6 Sol, दोन्ही Ultra effort.
% दुरुस्ती-पुनरावलोकन: GPT-6.1 Sol, Ultra effort.
\chapter{विन्यासमीमांसेचे अंकगणितीकरण}\label{inc:art::chap}
\begin{editorial}
  चिन्हांकित टॅब्लोसाठी अंकगणितीकरण अद्याप उपलब्ध नाही, याची नोंद घ्यावी.
\end{editorial}








\section{परिचय}\label{inc:art:int:sec}

संगणनक्षमता आणि तर्कशास्त्र यांचा संबंध जोडण्यासाठी तर्कशास्त्रातील
वस्तू (चिन्हे, पदे, सूत्रे, निष्पत्त्या), त्यांवरील क्रिया,
आणि त्यांचे गुणधर्म व संबंध यांविषयी संगणकीय हाताळणीस अनुकूल रीतीने
बोलता आले पाहिजे. चिन्हे, चिन्हांचे अनुक्रम आणि त्यांपासून तयार
होणाऱ्या इतर वस्तूंवरील संगणनक्षम फलने व संबंध विचारात घेऊन हे थेट
करता येते. तार्किक विन्यासमीमांसेतील सर्व वस्तू सांत असतात आणि
चिन्हांच्या गणनीय संचांपासून बनतात; म्हणून काही संगणन-प्रतिमानांत
अशी थेट हाताळणी शक्य आहे. पण पुनरावर्ती फलनांसारखी इतर संगणन-प्रतिमाने
संख्या आणि त्यांवरील संबंध व फलने यांपुरती मर्यादित असतात. शिवाय,
अखेरीस काही उपपत्त्यांच्या आतही विन्यासमीमांसा हाताळायची आहे;
विशेषतः अंकगणिताच्या भाषेत मांडलेल्या उपपत्त्यांमध्ये. अशा वेळी
विन्यासमीमांसेचे \emph{अंकगणितीकरण} करावे लागते: विन्यासात्मक वस्तू,
त्यांवरील क्रिया आणि त्यांचे संबंध यांचे अनुक्रमे संख्या, अंकगणितीय
फलने आणि अंकगणितीय संबंध यांच्या रूपात निरूपण करावे लागते. लाइब्निझपर्यंत
मागे जाणारी यामागील कल्पना म्हणजे विन्यासात्मक वस्तूंना संख्या नेमून देणे.

चिन्हांना त्यांचे ``संकेतांक'' म्हणून संख्या देणे तुलनेने सोपे आहे.
काही चिन्हांमुळे थोडी अडचण येते; उदाहरणार्थ, चले अनंत असतात
आणि प्रत्येक स्थानसंख्या~$n$ साठी फलने देखील अनंत असतात.
तरीही चिन्हांना पद्धतशीरपणे संख्या देता येतात, जेणेकरून, उदाहरणार्थ,
$\Obj v_2$ आणि $\Obj v_3$ यांना वेगवेगळे संकेतांक मिळतील. चिन्हांचे
अनुक्रम (उदा. पदे आणि सूत्रे) सांकेतिक करणे अधिक अवघड आहे.
पण संख्यांचे अनुक्रम निव्वळ अंकगणितीय रीतीने हाताळता आले (उदा.
अविभाज्य संख्यांच्या घातांनी अनुक्रम सांकेतिक करून), तर एकेका चिन्हाचे
सांकेतीकरण चिन्हांच्या अनुक्रमांपर्यंत वाढवता येते; त्यानंतर
सूत्रांचे अनुक्रम किंवा इतर मांडण्या, जसे निष्पत्त्या,
यांपर्यंतही ते वाढवता येते. या विस्तारित सांकेतीकरणाला ``गोडेल
क्रमांकन'' म्हणतात. प्रत्येक पदाला, सूत्राला आणि
निष्पत्ती ला गोडेल संख्या नेमली जाते.

चिन्हांचे अनुक्रम त्यांच्या संकेतांकांच्या अनुक्रमांनी सांकेतिक करून,
आणि संगणनक्षम फलनांनी हाताळता येईल अशी अनुक्रम-सांकेतीकरणाची पद्धत
निवडून, गोडेल संख्याही संगणनक्षम फलनांनी हाताळता येतात. प्रत्यक्षात
येथे लागणारी सर्व फलने आदिम पुनरावर्ती असतील. उदाहरणार्थ, अनुक्रमाच्या
संकेतांकावरून त्याची लांबी आणि त्यातील $i$-वा घटक काढण्याची दोन्ही
फलने आदिम पुनरावर्ती आहेत. अनुक्रम सांकेतिक करणारी संख्या, उदाहरणार्थ,
सूत्र~$\varphi$ ची गोडेल संख्या असेल, तर सूत्राची लांबी
आणि सूत्र मधील $i$-व्या चिन्हाचा संकेतांक हेही~$\varphi$ च्या
गोडेल संख्येवरून काढता येतात, हे लगेच दिसते. एखादी संख्या योग्य
रीतीने रचलेल्या पदाची किंवा योग्य निष्पत्तीची गोडेल संख्या
आहे हा गुणधर्म आदिम पुनरावर्ती आहे, हे सिद्ध करणे थोडे अधिक अवघड आहे.
तरीही ते शक्य आहे, कारण आपल्याला अभिप्रेत अनुक्रमांची (पदे,
सूत्रे, निष्पत्त्या) व्याख्या विगमनाधारित रीतीने केलेली आहे.

उदाहरण म्हणून आदेशनाची क्रिया विचारात घेऊ. $\varphi$ हे सूत्र, $x$ हा
चर आणि $t$ हे पद असेल, तर $\varphi$ मधील~$x$ च्या प्रत्येक मुक्त
आस्थितीऐवजी~$t$ ठेवल्यावर $\Subst{\varphi}{t}{x}$ मिळते. आता $\varphi$,
$x$, $t$ यांना अनुक्रमे $k$, $l$, $m$ या गोडेल संख्या दिल्या आहेत
असे समजू. तीच पद्धत $\Subst{\varphi}{t}{x}$ लाही एखादी गोडेल संख्या,
समजा~$n$, देते. $k$, $l$, $m$ यांना~$n$ शी जोडणारे हे फलन म्हणजे
आदेशन-क्रियेचे अंकगणितीय समरूप आहे. आदेशन-क्रिया $\varphi$, $x$, $t$
यांपासून $\Subst{\varphi}{t}{x}$ देते, तर अंकगणितीकृत आदेशन-फलन
$k$, $l$, $m$ या गोडेल संख्यांपासून गोडेल संख्या~$n$ देते.
हे फलन आदिम पुनरावर्ती आहे, हे आपण पाहू.

विन्यासमीमांसेचे अंकगणितीकरण केवळ अमूर्त रुचीचे नाही. अंकगणिताची
भाषा यांसारख्या भाषांत भाषांविषयी ``बोलण्याची'' साधने आधीपासून
नसूनही व्यंजकांचे गुंतागुंतीचे गुणधर्म त्यांत आकारिक रीतीने व्यक्त
करता येतात, ही कल्पना सुरुवातीला सहज सुचण्यासारखी नव्हती. मग
अशा भाषेतील, उदाहरणार्थ पेआनो अंकगणितासारखी, उपपत्ती स्वतःच्या
भाषेविषयी काय \emph{सिद्ध} करू शकते, हा प्रश्न विचारणे स्वाभाविक
ठरते (उदा. वाक्ये सिद्ध करता येतात का किंवा ती सत्य आहेत
का). यातून गोडेलची (असिद्धतायोग्यतेविषयी) आणि टार्स्कीची
(सत्याच्या अव्याख्येयतेविषयी) मर्यादा दाखवणारी प्रसिद्ध प्रमेये
समोर येतात. पण संगणनक्षमता उपपत्तीतील काही महत्त्वाचे निष्कर्ष
सिद्ध करण्यासाठीही विन्यासमीमांसेचे अंकगणितीकरण महत्त्वाचे आहे;
उदा. उपपत्त्यांची संगणकीय शक्ती किंवा संगणनक्षमतेच्या विविध
प्रतिमानांतील संबंध यांविषयीचे निष्कर्ष. इतर वस्तू आणि गुणधर्मांचे
अंकगणितीकरण कसे करावे याचा नमुनाही विन्यासमीमांसेचे अंकगणितीकरण
देते. उदाहरणार्थ, (ट्यूरिंग यंत्रांच्या) स्थितिवर्णनांचे आणि
संगणनांचेही अशाच रीतीने अंकगणितीकरण करता येते. त्यामुळे एका
प्रतिमानातील संगणनांचे (उदा. ट्यूरिंग यंत्रांतील) दुसऱ्या प्रतिमानात
(उदा. पुनरावर्ती फलनांत) अनुकरण करता येते.











\section{चिन्हांचे सांकेतीकरण}\label{inc:art:cod:sec}

प्रथम-क्रम तर्कशास्त्राच्या मूलभूत भाषा~$\Lang L$ मध्ये ही चिन्हे वापरली जातात:
\[\lfalse \quad \lnot \quad \lor \quad \land \quad \lif \quad \lforall
\quad \lexists \quad \eq \quad ( \quad ) \quad ,
\]
यांबरोबर चरे व स्थिरांक यांचे गणनीय संच, तसेच कोणत्याही
स्थानसंख्येची फलने व विधेये यांचे गणनीय संचही
वापरले जातात. या प्रत्येक चिन्हाला \emph{संकेतांक} म्हणून एक संख्या
देता येते: प्रत्येक चिन्हाचा संकेतांक वेगळा असेल आणि दोन भिन्न
चिन्हांना एकच संख्या दिली जाणार नाही. सर्व चिन्हांचा संच गणनीय
असल्यामुळे त्याची नैसर्गिक संख्यांच्या संचाशी एकास-एक आच्छादन आहे;
म्हणून हे शक्य आहे. परंतु संकेतांकावरून चिन्ह आणि त्याविषयीची काही
माहिती (उदा. फलन ची स्थानसंख्या) संगणनक्षम रीतीने परत
मिळाली पाहिजे. असे करण्याचे अनेक मार्ग आहेत. पुढील मार्गात आदिम
पुनरावर्ती फलने वापरली आहेत. ($\tuple{n_0, \dots, n_k}$ ही
$n_0$, \dots, $n_k$ या संख्यांच्या अनुक्रमाला सांकेतिक करणारी संख्या
आहे, हे आठवा.)

\begin{defn}
$s$ हे~$\Lang L$ मधील चिन्ह असेल, तर त्याचा \emph{चिन्हसंकेतांक}~$\scode s$
पुढीलप्रमाणे परिभाषित करू:
\begin{enumerate}
\item $s$ हे तार्किक चिन्हांपैकी असेल, तर पुढील सारणीनुसार $\scode s$ ठरतो:
\[\begin{array}{cccccccccc}
  \lfalse & \lnot & \lor & \land & \lif & \lforall \\
\tuple{0, 0} & \tuple{0, 1} & \tuple{0, 2} & \tuple{0, 3} &
\tuple{0, 4} & \tuple{0, 5} \\
\lexists & \eq & ( & ) & ,\\
\tuple{0, 6} & \tuple{0, 7} &
\tuple{0, 8} & \tuple{0, 9} & \tuple{0, 10}
\end{array}
\]
\item $s$ हा $i$-वा चर $\Obj v_i$ असेल, तर $\scode s = \tuple{1, i}$.
\item $s$ हे $i$-वे स्थिरांक~$\Obj c_i$ असेल, तर
  $\scode s = \tuple{2, i}$.
\item $s$ हे $i$-वे $n$-स्थानी फलन~$\Obj f_i^n$ असेल, तर
  $\scode s = \tuple{3, n, i}$.
\item $s$ हे $i$-वे $n$-स्थानी विधेय~$\Obj P_i^n$ असेल, तर
  $\scode s = \tuple{4, n, i}$.
\end{enumerate}
\end{defn}

\begin{prop}
पुढील संबंध आदिम पुनरावर्ती आहेत:
\begin{enumerate}
\item $\fn{Fn}(x, n)$ तेव्हा आणि केवळ तेव्हाच जेव्हा $x$ हा एखाद्या~$i$ साठी
  $\Obj f^n_i$ चा संकेतांक असतो; म्हणजेच $x$ हा $n$-स्थानी
  फलन चा संकेतांक असतो.
\item $\fn{Pred}(x, n)$ तेव्हा आणि केवळ तेव्हाच जेव्हा $x$ हा एखाद्या~$i$ साठी
  $\Obj P^n_i$ चा संकेतांक असतो, किंवा $x$ हा $\eq$ चा संकेतांक
  असतो आणि $n = 2$ असते; म्हणजेच $x$ हा $n$-स्थानी विधेय
  चा संकेतांक असतो.
\end{enumerate}
\end{prop}

\begin{defn}
$s_0, \dots, s_{n-1}$ हा चिन्हांचा अनुक्रम असेल, तर त्याची
\emph{गोडेल संख्या} $\tuple{\scode{s_0}, \dots, \scode{s_{n-1}}}$ आहे.
\end{defn}

\begin{explain}
\emph{संकेतांक} आणि \emph{गोडेल संख्या} वेगवेगळ्या गोष्टी आहेत,
हे लक्षात घ्या. उदाहरणार्थ, $\Obj v_5$ या चराचा संकेतांक
$\scode{\Obj v_5} = \tuple{1, 5} = 2^2\cdot 3^6$ आहे.
पण पद म्हणून पाहिले असता~$\Obj v_5$ हाही लांबी~$1$ असलेला
चिन्हांचा अनुक्रम आहे. त्यामुळे $\Obj v_5$ या \emph{पदाची}
\emph{गोडेल संख्या}~$\Gn{\Obj v_5}$ ही
$\tuple{\scode{\Obj v_5}} = 2^{\scode{\Obj v_5} + 1} =
2^{2^2\cdot 3^6 + 1}$ आहे.
\end{explain}

\begin{ex}
$k_0$, \dots, $k_{n-1}$ हा संख्यांचा अनुक्रम असेल, तर
अविभाज्य संख्यांच्या घातांनी केलेल्या सांकेतीकरणात
$\tuple{k_0, \dots, k_{n-1}}$ हा त्याचा संकेतांक
\[2^{k_0+1}\cdot3^{k_1+1}\cdot \dots \cdot p_{n-1}^{k_{n-1}+1},
\]
आहे, हे आठवा. येथे $p_i$ ही $i$-वी अविभाज्य संख्या आहे
($p_0 = 2$ पासून सुरुवात होते). म्हणून, उदाहरणार्थ,
$\eq[\Obj v_0][\Obj 0]$ हे सूत्र, किंवा अधिक स्पष्टपणे
${\eq}(\Obj v_0,\Obj c_0)$, याची गोडेल संख्या
\[\tuple{\scode{\eq},\scode{(},\scode{\Obj v_0},\scode{,}, \scode{\Obj
    c_0},\scode{)}}.
\]
आहे. येथे $\scode{\eq}$ हा $\tuple{0,7} = 2^{0+1}\cdot
3^{7+1}$ आहे; $\scode{\Obj v_0}$ हा $\tuple{1,0} = 2^{1+1}\cdot3^{0+1}$
आहे; इत्यादी. म्हणून $\Gn{=(\Obj v_0,\Obj c_0)}$ ही संख्या
\begin{multline*}
2^{\scode{=} + 1}\cdot 3^{\scode{(}+1}\cdot 5^{\scode{\Obj v_0}+1}
\cdot 7^{\scode{,} + 1} \cdot 11^{\scode{\Obj c_0}+1} \cdot
13^{\scode{)}+1} = \\
2^{2^1\cdot 3^8 + 1}\cdot 3^{2^1\cdot 3^9+1}\cdot 5^{2^2\cdot 3^1+1}
\cdot 7^{2^1\cdot 3^{11} + 1} \cdot 11^{2^3\cdot3^1+1} \cdot
13^{2^1\cdot3^{10}+1} = \\
2^{13\,123}\cdot 3^{39\,367}\cdot 5^{13}\cdot 7^{354\,295}\cdot11^{25}\cdot13^{118\,099}.
\end{multline*}
\end{ex}











\section{पदांचे सांकेतीकरण}\label{inc:art:trm:sec}

\begin{explain}
पद हा चिन्हांच्या अनुक्रमाचा एक विशिष्ट प्रकार आहे: पदांसाठीच्या
रचनानियमांनुसार स्थिरचिन्हे आणि चरे यांपासून त्याची विगमनाधारित
रचना होते. चिन्हे सांकेतिक करण्याची पद्धत आणि संख्यांचे अनुक्रम
सांकेतिक करण्याची पद्धत वापरून चिन्हांचे अनुक्रम संख्यांनी सांकेतिक
करता येतात. त्यामुळे पदांना गोडेल संख्या देणे कठीण नाही. खरी
अडचण ही आहे की एखादी संख्या योग्य रीतीने रचलेल्या पदाची गोडेल
संख्या आहे हा गुणधर्म संगणनक्षम, किंबहुना आदिम पुनरावर्ती आहे, हे दाखवणे.
\end{explain}

चले आणि स्थिरांक ही सर्वांत सोपी पदे आहेत. $x$ ही
अशा पदाची गोडेल संख्या आहे का, हे तपासणे सोपे आहे:
$x$ ही एखाद्या~$i$ साठी $\Gn{\Obj v_i}$ असेल, तर आणि तरच
$\fn{Var}(x)$ खरे असते. दुसऱ्या शब्दांत, $x$ हा लांबी~$1$
असलेल्या अनुक्रमाचा संकेतांक आहे आणि त्यातील एकमेव घटक $(x)_0$
हा एखाद्या चल~$\Obj v_i$ चा संकेतांक आहे; म्हणजे एखाद्या~$i$
साठी $x$ हा $\tuple{\tuple{1, i}}$ आहे. तसेच, $x$ ही एखाद्या~$i$
साठी $\Gn{\Obj c_i}$ असेल, तर आणि तरच $\fn{Const}(x)$ खरे असते.
हे दोन्ही संबंध आदिम पुनरावर्ती आहेत; कारण असा~$i$ अस्तित्वात
असेल, तर तो $< x$ असलाच पाहिजे:
\begin{align*}
  \fn{Var}(x) & \defiff \bexists{i<x}{x = \tuple{\tuple{1, i}}}\\
  \fn{Const}(x) & \defiff \bexists{i<x}{x = \tuple{\tuple{2, i}}}
\end{align*}

\begin{prop}
\label{inc:art:trm:prop:term-primrec}
$x$ ही अनुक्रमे पदाची किंवा बंद पदाची गोडेल संख्या असेल, तर आणि तरच
खरे असणारे $\fn{Term}(x)$ आणि $\fn{ClTerm}(x)$ हे संबंध आदिम पुनरावर्ती आहेत.
\end{prop}

\begin{proof}
चिन्हांचा अनुक्रम~$s$ हा पद असेल, तर आणि तरच पदांची अशी
रचनाक्रमिका~$s_0$, \dots, $s_{k-1} = s$ अस्तित्वात असते की जिच्यात
पदांच्या रचनानियमांनुसार स्थिरांक आणि चले यांपासून
पद~$s$ कसे बनले याची नोंद असते. एखादी प्रस्तावित रचनाक्रमिका
रचनानियमांचे पालन करते ही अट व्यक्त करण्यासाठी, प्रत्येक $i < k$
बाबत पुढीलपैकी एक अट खरी असली पाहिजे:
\begin{enumerate}
\item $s_i$ हे चल~$\Obj v_j$ आहे; किंवा
\item $s_i$ हे स्थिरांक~$\Obj c_j$ आहे; किंवा
\item $i$-व्या स्थानापूर्वी आलेल्या $t_1$, \dots, $t_n$ या $n$
  पदांपासून $n$-स्थानी फलन~$\Obj f^n_j$ वापरून $s_i$ बनले आहे.
\end{enumerate}
गोडेल संख्यांवरील संबंधित संबंध आदिम पुनरावर्ती आहे हे दाखवण्यासाठी
ही अट आदिम पुनरावर्ती फलने, संबंध आणि परिबद्ध संख्यकीकरण वापरून
व्यक्त करावी लागेल.

$y$ ही $s_0$, \dots, $s_{k-1}$ या अनुक्रमाचा संकेतांक आहे असे समजू;
म्हणजे $y = \tuple{\Gn{s_0}, \dots, \Gn{s_{k-1}}}$. तो गोडेल
संख्या~$x$ असलेल्या पदाच्या रचनाक्रमिकेचा संकेतांक असेल, तर आणि तरच
प्रत्येक $i < k$ साठी पुढीलपैकी एक अट खरी असते:
\begin{enumerate}
\item $\fn{Var}((y)_i)$; किंवा
\item $\fn{Const}((y)_i)$; किंवा
\item एखादा~$n$ आणि $z = \tuple{z_1, \dots, z_n}$ ही संख्या
  अस्तित्वात असतात, अशी की प्रत्येक $z_l$ हा $i' < i$ असलेल्या
  एखाद्या निर्देशांकासाठी $(y)_{i'}$ बरोबर असतो आणि
\[(y)_i = \Gn{\Obj f^n_j(} \concat \fn{flatten}(z) \concat \Gn{)},
\]
\end{enumerate}
आणि याशिवाय $(y)_{k-1} = x$ असते. ($\fn{flatten}(z)$ हे फलन
$\tuple{\Gn{t_1}, \dots, \Gn{t_n}}$ या अनुक्रमापासून
$\Gn{t_1, \dots, t_n}$ तयार करते आणि ते आदिम पुनरावर्ती आहे.)

(3) मधील $j$, $n$ हे निर्देशांक, $t_l$ या पदांच्या गोडेल संख्या
$z_l$, आणि $\tuple{z_1, \dots, z_n}$ या अनुक्रमाचा संकेतांक~$z$
हे सर्व~$y$ पेक्षा लहान आहेत. वरच्या~$k$ ऐवजी $\len{y}$
वापरता येते. म्हणून ``$y$ ही गोडेल संख्या~$x$ असलेल्या पदाच्या
रचनाक्रमिकेचा संकेतांक आहे'' हे वाक्य त्या संबंधाचे आदिम पुनरावर्ती
असणे स्पष्ट करणाऱ्या रीतीने व्यक्त करता येते.

आता~$y$ साठी आदिम पुनरावर्ती मर्यादा आहे याचीच खात्री करून घ्यायची
आहे. $x$ ही एखाद्या पदाची गोडेल संख्या असेल, तर त्याला जास्तीत
जास्त $\len{x}$ पदांची रचनाक्रमिका असली पाहिजे: कारण~$s$ च्या
केवळ उपपदांची रचनाक्रमिका निवडता येते. तिच्यातील प्रत्येक पदाची
सुरुवात~$s$ मधील एखाद्या स्थानावर
होते, आणि दोन उपपदे एकाच स्थानावर सुरू होऊ शकत नाहीत. $s$ च्या
प्रत्येक उपपदाची गोडेल संख्या अर्थातच $\le x$ आहे. म्हणून
$k=\len{x}$ असेल, तर संकेतांक $\le p_{k-1}^{k(x+1)}$ असलेली
रचनाक्रमिका नेहमी अस्तित्वात असते.


$\fn{ClTerm}$ साठी चले संबंधी अट फक्त वगळावी.
\end{proof}

\begin{prob}
$\tuple{\Gn{t_1}, \dots, \Gn{t_n}}$ या अनुक्रमापासून
$\Gn{t_1, \dots, t_n}$ तयार करणारे $\fn{flatten}(z)$ हे फलन
आदिम पुनरावर्ती आहे, हे दाखवा.
\end{prob}

\begin{prop}\label{inc:art:trm:prop:num-primrec}
  $\fn{num}(n) = \Gn{\num{n}}$ हे फलन आदिम पुनरावर्ती आहे.
\end{prop}

\begin{proof}
  $\fn{num}(n)$ ची आदिम पुनरावर्तनाने व्याख्या करू:
  \begin{align*}
    \fn{num}(0) & = \Gn{\Obj 0}\\
    \fn{num}(n+1) & = \Gn{\prime(} \concat \fn{num}(n) \concat \Gn{)}.
  \end{align*}
\end{proof}











\section{सूत्र यांचे सांकेतीकरण}\label{inc:art:frm:sec}

$\fn{Term}(x)$ हा संबंध आदिम पुनरावर्ती रीतीने परिभाषित केल्यावर
सूत्रे साठीचा संबंधित संबंध $\fn{Frm}(x)$ हाही आदिम
पुनरावर्ती रीतीने परिभाषित करता येतो.

\begin{prop}
$x$ ही आण्विक सूत्राची गोडेल संख्या असेल, तर आणि तरच
खरे असणारा $\fn{Atom}(x)$ हा संबंध आदिम पुनरावर्ती आहे.
\end{prop}

\begin{proof}
$x$ ही आण्विक सूत्राची गोडेल संख्या असेल, तर आणि तरच
पुढीलपैकी एक अट खरी असते:
\begin{enumerate}
\item $n$, $j < x$ आणि $z \le B(x)$ अशी मूल्ये अस्तित्वात आहेत की
  $\len{z}=n$ आणि प्रत्येक $i < n$ साठी $\fn{Term}((z)_i)$ खरे आहे आणि $x = $
\[\Gn{\Obj P^n_j(} \concat \fn{flatten}(z) \concat \Gn{)}.
\]
\item $z_1, z_2 < x$ अशी मूल्ये अस्तित्वात आहेत की $\fn{Term}(z_1)$,
  $\fn{Term}(z_2)$ खरे आहेत आणि $x = $
\[\Gn{{\eq}(} \concat z_1 \concat \Gn{,} \concat z_2 \concat{\Gn{)}}.
\]
  \item $x = \Gn{\lfalse}$.
  \item $x = \Gn{\ltrue}$.
\end{enumerate}
येथे $k=\len{x}\ge1$ आणि $B(x)=p_{k-1}^{k(x+1)}$ घ्या.
शून्य लांबीचा अनुक्रम आण्विक सूत्र नाही. प्रत्येक आदानपदाचा संकेतांक
अंतिम सूत्राच्या संकेतांकापेक्षा लहान असतो; मात्र त्या संकेतांकांच्या
संपूर्ण अनुक्रमाचा संकेतांक तसाच लहान असेल असे नाही. आदानपदांची
संख्या $n\le k$ असल्याने $B(x)$ ही त्या अनुक्रमासाठी आदिम पुनरावर्ती
वरची मर्यादा आहे.

\end{proof}

\begin{prop}
\label{inc:art:frm:prop:frm-primrec}
$x$ ही सूत्राची गोडेल संख्या असेल, तर आणि तरच
खरे असणारा $\fn{Frm}(x)$ हा संबंध आदिम पुनरावर्ती आहे.
\end{prop}

\begin{proof}
चिन्हांचा अनुक्रम~$s$ हे सूत्र असेल, तर आणि तरच
सूत्राची अशी रचनाक्रमिका~$s_0$, \dots, $s_{k-1} = s$
अस्तित्वात असते की जिच्यात रचनानियमांनुसार आण्विक सूत्रे
पासून~$s$ कसे बनले याची नोंद असते. केवळ उपसूत्रांची रचनाक्रमिका
निवडल्यास प्रत्येक $s_i$ चा संकेतांक
$s$ च्या संकेतांक~$x$ पेक्षा मोठा नसतो; पण संपूर्ण रचनाक्रमिका
$\tuple{s_0, \dots, s_{k-1}}$ हिचा संकेतांक त्यापेक्षा लहान
असणे आवश्यक नाही: मूळ इंग्रजीतील उलट दावा एका आण्विक सूत्रासाठीही
खोटा ठरतो. त्याऐवजी मागील पदांच्या सिद्धतेप्रमाणे त्याला
त्या अंतिम सूत्राच्या संकेतांकावर अवलंबून आदिम पुनरावर्ती वरची
मर्यादा देता येते.
\end{proof}

\begin{prob}
\ref{inc:art:frm:prop:frm-primrec} ची,
\ref{inc:art:trm:prop:term-primrec} च्या पहिल्या सिद्धतेच्या
धर्तीवर, सविस्तर सिद्धता द्या.
\end{prob}

\begin{prop}
\label{inc:art:frm:prop:freeocc-primrec}
$x$ ही गोडेल संख्या असलेल्या सूत्रातील $i$-वे चिन्ह हे
गोडेल संख्या~$z$ असलेल्या चराची मुक्त आस्थिती असेल,
तर आणि तरच खरे असणारा $\fn{FreeOcc}(x, z, i)$ हा संबंध
आदिम पुनरावर्ती आहे.
\end{prop}

\begin{proof}
सराव.
\end{proof}

\begin{prob}
\ref{inc:art:frm:prop:freeocc-primrec} सिद्ध करा.
सूत्राच्या ज्या उपचिन्हमाला स्वतः सूत्र आहेत,
त्या त्याची उप-सूत्र आहेत, हे तथ्य वापरू शकता.
\end{prob}

\begin{prop}
$x$ ही वाक्याची गोडेल संख्या असेल, तर आणि तरच
खरे असणारा $\fn{Sent}(x)$ हा गुणधर्म आदिम पुनरावर्ती आहे.
\end{prop}

\begin{proof}
वाक्य म्हणजे चलांच्या मुक्त आस्थिती नसलेले
सूत्र. म्हणून $\fn{Sent}(x)$ तेव्हा आणि केवळ तेव्हाच खरे असते जेव्हा
\begin{multline*}
\fn{Frm}(x) \land \bforall{i<\len{x}}{\bforall{z<x}{}}\\
(\bexists{j<z}{z=\Gn{\Obj v_j}} \lif \lnot\fn{FreeOcc}(x,z,i)).
\end{multline*}

\end{proof}












\section{आदेशन}\label{inc:art:sub:sec}

आदेशन म्हणजे सूत्र~$\varphi$ मध्ये चल~$u$ च्या
सर्व मुक्त आस्थितींच्या जागी पद~$t$ ठेवण्याची क्रिया आहे;
ती $\Subst{\varphi}{t}{u}$ अशी लिहितात, हे आठवा. चले,
सूत्रे आणि पदे यांच्या गोडेल संख्यांवर ही क्रिया
केली, तर ती आदिम पुनरावर्ती असते.

\begin{prop}\label{inc:art:sub:prop:subst-primrec}
पुढील गुणधर्म असलेले आदिम पुनरावर्ती फलन $\fn{Subst}(x, y, z)$
अस्तित्वात आहे:
\[\fn{Subst}(\Gn{\varphi}, \Gn{t}, \Gn{u}) = \Gn{\Subst{\varphi}{t}{u}}.
\]
\end{prop}

\begin{proof}
आता $\fn{hSubst}$ हे फलन पुढीलप्रमाणे आदिम पुनरावर्तनाने
परिभाषित करता येते:
\begin{multline*}
\begin{aligned}
\fn{hSubst}(x, y, z, 0) & = \emptyseq \\
\fn{hSubst}(x, y, z, i+1) & =
\end{aligned}\\
\begin{cases}
\fn{hSubst}(x, y, z, i) \concat y & \text{जर $\fn{FreeOcc}(x, z, i)$ खरे असेल तर} \\
\fn{append}(\fn{hSubst}(x, y, z, i), (x)_{i}) & \text{अन्यथा.}
\end{cases}
\end{multline*}
आता $\fn{Subst}(x, y, z)$ ची व्याख्या
$\fn{hSubst}(x, y, z, \len{x})$ अशी करता येते.
\end{proof}

\begin{prop}
\label{inc:art:sub:prop:free-for}
गोडेल संख्या~$y$ असलेले पद, गोडेल संख्या~$x$ असलेल्या सूत्रात
गोडेल संख्या~$z$ असलेल्या चरासाठी आदेशनासाठी मुक्त असेल,
तर आणि तरच खरे असणारा $\fn{FreeFor}(x, y, z)$ हा संबंध
आदिम पुनरावर्ती आहे.
\end{prop}

\begin{proof} सराव. \end{proof}

\begin{prob}
\ref{inc:art:sub:prop:free-for} सिद्ध करा.
\end{prob}











\section{$\Log{LK}$ मधील निष्पत्ती}\label{inc:art:plk:sec}

\begin{explain}
निष्पत्त्यांचे अंकगणितीकरण करण्यासाठी निष्पत्त्या संख्यांनी दर्शवावे लागतात.
प्रत्येक निगमनाला नामांकन असलेले क्रमवर्तींचे वृक्ष म्हणजे निष्पत्त्या
असल्यामुळे, पुनरावर्ती निरूपण हा सहज मार्ग आहे. आपण निष्पत्ती
एका क्रमित समूहाने दर्शवतो: त्याचे घटक म्हणजे अंतिम क्रमवर्ती, नामांकन,
आणि अखेरच्या निगमनाच्या आधारविधानांपर्यंत येणाऱ्या उप-निष्पत्त्यांची निरूपणे.
\end{explain}

\begin{defn}
$\Gamma$ ही वाक्यांची सांत क्रमिका, $\Gamma =
\tuple{\varphi_1, \dots, \varphi_n}$, असेल, तर $\Gn{\Gamma} = \tuple{\Gn{\varphi_1},
  \dots, \Gn{\varphi_n}}$.

$\Gamma \Sequent \Delta$ ही क्रमवर्ती असेल, तर
$\Gamma \Sequent \Delta$ ची गोडेल संख्या
पुढीलप्रमाणे असते:
\[\Gn{\Gamma \Sequent \Delta} = \tuple{\Gn{\Gamma}, \Gn{\Delta}}
\]

$\pi$ ही $\Log{LK}$ मधील निष्पत्ती असेल, तर $\Gn{\pi}$ ची
व्याख्या पुढीलप्रमाणे करतो:
\begin{enumerate}
\item $\pi$ मध्ये फक्त $\Gamma \Sequent \Delta$ ही आरंभीची क्रमवर्ती
  असेल, तर $\Gn{\pi}$ म्हणजे
  \[    \tuple{0, \Gn{\Gamma \Sequent \Delta}}.
  \]
\item $\pi$ चा शेवट एक किंवा दोन आधारविधानांच्या निगमनाने होत असेल,
  त्या निगमनाचा निष्कर्ष $\Gamma \Sequent \Delta$ असेल, आणि $\pi_1$
  व $\pi_2$ या अखेरच्या निगमनाच्या आधारविधानांवर संपणाऱ्या तात्काळ
  उपनिष्पत्ती असतील, तर $\Gn{\pi}$ अनुक्रमे
  \begin{align*}
    & \tuple{1, \Gn{\pi_1}, \Gn{\Gamma \Sequent \Delta}, k} \text{ किंवा}\\
    & \tuple{2, \Gn{\pi_1}, \Gn{\pi_2}, \Gn{\Gamma \Sequent \Delta}, k},
  \end{align*}
  असते. अखेरच्या निगमनात कोणता नियम वापरला आहे त्यानुसार $k$ पुढील
  तक्त्यात दिले आहे:

  \begin{tabular}{lccccccc}
    \text{नियम:} & \LeftR{\Weakening} & \RightR{\Weakening} &
    \LeftR{\Contraction} & \RightR{\Contraction} &
    \LeftR{\Exchange} & \RightR{\Exchange} \\
    $k$: & 1 & 2 & 3 & 4 & 5 & 6 \\[2ex]
    \text{नियम:} & \LeftR{\lnot} & \RightR{\lnot} &
    \LeftR{\land} & \RightR{\land} &
    \LeftR{\lor} & \RightR{\lor} \\
    $k$: & 7 & 8 & 9 & 10 & 11 & 12 \\[2ex]
    \text{नियम:} & \LeftR{\lif} & \RightR{\lif} &
    \LeftR{\lforall} & \RightR{\lforall} &
    \LeftR{\lexists} & \RightR{\lexists} \\
    $k$: & 13 & 14 & 15 & 16 & 17 & 18 \\[2ex]
    \text{नियम:} & \Cut & = \\
    $k$: & 19 & 20
  \end{tabular}
\end{enumerate}
\end{defn}
\begin{ex}
  पुढील अगदी सोपी निष्पत्ती पाहा:
  \begin{prooftree}
    \Axiom$\varphi \fCenter \varphi$
    \RightLabel{\LeftR{\land}}
    \UnaryInf$\varphi \land \psi \fCenter \varphi$
    \RightLabel{\RightR{\lif}}
    \UnaryInf$\fCenter (\varphi \land \psi) \lif \varphi$
  \end{prooftree}
  आरंभीच्या क्रमवर्तीची गोडेल संख्या $p_0 = \tuple{0,
    \Gn{\varphi \Sequent \varphi}}$ असेल. $\LeftR{\land}$ च्या निष्कर्षावर
    संपणाऱ्या निष्पत्तीची गोडेल संख्या $p_1 =
    \tuple{1, p_0, \Gn{\varphi \land \psi \Sequent \varphi}, 9}$ असेल. येथे $1$
    हे $\LeftR{\land}$ ला एकच आधारविधान असल्यामुळे आले आहे;
    $\varphi \land \psi \Sequent \varphi$ या निष्कर्षाची गोडेल संख्या पुढे आहे;
    आणि $9$ हा $\LeftR{\land}$ चा संकेतांक आहे. मग संपूर्ण
    निष्पत्तीची गोडेल संख्या
    $\tuple{1, p_1, \Gn{\Sequent (\varphi \land \psi) \lif \varphi}, 14}$,
    म्हणजेच
  \[
  \tuple{1, \tuple{1, \tuple{0, \Gn{\varphi \Sequent \varphi}}, \Gn{\varphi \land \psi \Sequent \varphi}, 9},
    \Gn{\Sequent (\varphi \land \psi) \lif \varphi}, 14}.
  \]
\end{ex}

\begin{explain}
निष्पत्त्यांचे निरूपण ठरल्यानंतर, अशा निष्पत्त्यांवर आदिम पुनरावर्ती रीतीने
क्रिया करता येतात आणि त्यांचे आवश्यक गुणधर्म व संबंध तसेच व्यक्त करता
येतात, हेही दाखवावे लागेल. काही क्रिया सोप्या आहेत. उदाहरणार्थ,
निष्पत्तीची गोडेल संख्या~$p$ दिली असता,
$\fn{EndSequent}(p) = (p)_{(p)_0+1}$ हे तिच्या अंतिम क्रमवर्तीची
गोडेल संख्या देते आणि $\fn{LastRule}(p) = (p)_{(p)_0+2}$ हे तिच्या
अखेरच्या नियमाचा संकेतांक देते. पुढीलप्रमाणे परिभाषित केलेला
क्रमिकेच्या संकेतांकाची वैधता तपासण्यासाठी आदिम पुनरावर्ती अट घ्या:
\[\fn{Seq}(u) \defiff u=0 \lor u=\prod_{i<\len{u}}p_i^{(u)_i+1}.
\]
ती रिक्त क्रमिकेचे दोन्ही संकेतांक स्वीकारते आणि अतिरिक्त अविभाज्य
घटक असलेले अवैध संकेतांक वगळते.
$\fn{Sequent}(s)$ गुणधर्म
\[  \begin{gathered}
  s=\tuple{(s)_0,(s)_1} \land \fn{Seq}((s)_0) \land \fn{Seq}((s)_1)\\
  {}\land \bforall{i<\len{(s)_0} + \len{(s)_1}}{\fn{Sent}(((s)_0 \concat (s)_1)_i)}
  \end{gathered}
\]
$s$ साठी खरा असतो, तर आणि तरच $s$ ही वाक्ये असलेल्या
क्रमवर्तीची गोडेल संख्या असते. काही इतर क्रिया अधिक कठीण आहेत; त्या कशा कराव्यात याची
किमान रूपरेषा आपण पाहू. $\pi$ ही $\Gamma$ पासून $\varphi$ ची
निष्पत्ती आहे, हा संबंध $\pi$ आणि $\varphi$ यांच्या गोडेल संख्यांचा
आदिम पुनरावर्ती संबंध आहे, हे दाखवणे आपले उद्दिष्ट आहे.
\end{explain}

\begin{prop}\label{inc:art:plk:prop:followsby}
  गोडेल संख्या~$p$ असलेल्या निष्पत्ती~$\pi$ मधील अखेरचे निगमन
  योग्य असेल, तर आणि तरच खरा असणारा $\fn{Correct}(p)$ गुणधर्म
  आदिम पुनरावर्ती आहे.
\end{prop}

\begin{proof}
  $\Gamma \Sequent \Delta$ ही आरंभीची क्रमवर्ती दोनपैकी एका
  परिस्थितीत असते: एखादे वाक्य~$\varphi$ असे असते की
  $\Gamma \Sequent \Delta$ म्हणजे $\varphi \Sequent \varphi$ असते;
  किंवा एखादे बंद पद~$t$ असे असते की
  $\Gamma \Sequent \Delta$ म्हणजे $\emptyset \Sequent \eq[t][t]$ असते.
  गोडेल संख्यांच्या भाषेत, मूलभूत आरंभीच्या क्रमवर्तींसाठी
  $\fn{InitialSeq}(s)$ पुढील अट पूर्ण झाली तर आणि तरच खरे असते;
  चिन्हांकित आवृत्त्यांत खाली दिलेल्या अटीही जोडायच्या आहेत:
  \begin{align*}
  \bexists{x < s}{} (\fn{Sent}(x) & \land
  s = \tuple{\tuple{x},\tuple{x}})
  \lor {}\\
  \bexists{t<s}{} (\fn{ClTerm}(t) & \land
  s = \tuple{0, \tuple{\Gn{{\eq}(} \concat t \concat \Gn{,} \concat t \concat \Gn{)}}}).
  \end{align*}
  या चिन्हांकित आवृत्तीत $s=\tuple{0,\tuple{\Gn{\ltrue}}}$
  हीदेखील आरंभीची क्रमवर्ती आहे.
  या चिन्हांकित आवृत्तीत $s=\tuple{\tuple{\Gn{\lfalse}},0}$
  हीदेखील आरंभीची क्रमवर्ती आहे.

  प्रत्येक निगमन नियम~$R$ साठी $\fn{FollowsBy}_R(p)$ हा संबंधही
  आदिम पुनरावर्ती आहे, हे दाखवावे लागेल. $p$ ही निष्पत्ती~$\pi$
  ची गोडेल संख्या असेल आणि $\pi$ ची अंतिम क्रमवर्ती तिच्या तात्काळ
  $\pi$ च्या उप-निष्पत्त्यांच्या अंतिम क्रमवर्तींपासून $R$ च्या योग्य वापराने
  निष्पन्न होत असेल, तर आणि तरच $\fn{FollowsBy}_R(p)$ खरे असते.

  \RightR{\land} नियमाचे उदाहरण सोपे आहे. $\pi$ चा शेवट
  $\RightR{\land}$ च्या योग्य वापराने होत असेल, तर तिचे रूप असे असते:
  \begin{prooftree}
    \AxiomC{}
    \RightLabel{$\pi_1$}
    \Deduce$\Gamma \fCenter \Delta, \varphi$

    \AxiomC{}
    \RightLabel{$\pi_2$}
    \Deduce$\Gamma \fCenter \Delta, \psi$

    \RightLabel{\RightR\land}
    \BinaryInf$\Gamma \fCenter \Delta, \varphi \land \psi$
  \end{prooftree}
  त्यामुळे निष्पत्ती~$\pi$ मधील अखेरचे निगमन $\RightR{\land}$ चा योग्य वापर
  असेल, तर आणि तरच $\Gamma$ व $\Delta$ या वाक्यांच्या क्रमिका
  आणि $\varphi$, $\psi$ ही दोन वाक्ये अशी असतात की $\pi_1$ ची
  अंतिम क्रमवर्ती $\Gamma \Sequent \Delta, \varphi$, $\pi_2$ ची अंतिम
  क्रमवर्ती $\Gamma \Sequent \Delta, \psi$, आणि $\pi$ ची अंतिम
  क्रमवर्ती $\Gamma \Sequent \Delta, \varphi \land \psi$ असते.
  हेच आता गोडेल संख्यांत मांडायचे आहे. $s = \Gn{\Gamma \Sequent \Delta}$
  असेल, तर $(s)_0 = \Gn{\Gamma}$ आणि $(s)_1 = \Gn{\Delta}$.
  म्हणून $\fn{FollowsBy}_{\RightR{\land}}(p)$ पुढील अट पूर्ण झाली
  तर आणि तरच खरे असते:
\begin{align*}
& \bexists{g < p}{\bexists{d < p}{\bexists{a < p}{\bexists{b < p}{\quad}}}} \\
& \qquad \fn{Sequent}(\tuple{g,d}) \land \fn{Sent}(a) \land \fn{Sent}(b) \land {}\\
& \qquad \fn{EndSequent}(p) =
  \tuple{g, d \concat
    \tuple{\Gn{(} \concat a \concat \Gn{\land} \concat b \concat \Gn{)}}} \land {} \\
& \qquad \fn{EndSequent}((p)_1) = \tuple{g, d \concat \tuple{a}} \land {}\\
& \qquad \fn{EndSequent}((p)_2) = \tuple{g, d \concat \tuple{b}} \land {}\\
& \qquad (p)_0 = 2 \land \fn{LastRule}(p) = 10.
\end{align*}
या ओळी अनुक्रमे पुढील गोष्टी व्यक्त करतात: $\Gamma$ ची गोडेल संख्या
$g$, $\Delta$ ची गोडेल संख्या $d$, सूत्र~$\varphi$ ची गोडेल संख्या $a$, आणि
सूत्र~$\psi$ ची गोडेल संख्या $b$ आहे; $\pi$ ची अंतिम क्रमवर्ती
$\Gamma \Sequent \Delta, \varphi \land \psi$ आहे; $\pi_1$ ची अंतिम क्रमवर्ती
$\Gamma \Sequent \Delta, \varphi$ आहे; $\pi_2$ ची अंतिम क्रमवर्ती
$\Gamma \Sequent \Delta, \psi$ आहे; आणि $\pi$ ला दोन तात्काळ
उपनिष्पत्ती असून तिचा अखेरचा नियम क्रमांक~$10$ असलेला
$\RightR\land$ आहे.

$\pi$ मधील अखेरचे निगमन $\RightR\lexists$ चा योग्य वापर असेल,
तर आणि तरच $\Gamma$ व $\Delta$ या क्रमिका, $\varphi$ हे सूत्र,
$x$ हे चल आणि $t$ हे बंद पद अशी असतात की $\pi$ ची अंतिम क्रमवर्ती
$\Gamma \Sequent \Delta, \lexists[x][\varphi]$ आणि $\pi_1$ ची अंतिम
क्रमवर्ती $\Gamma \Sequent \Delta, \Subst{\varphi}{t}{x}$ असते.
म्हणून गोडेल संख्यांत $\fn{FollowsBy}_{\RightR\lexists}(p)$ पुढील
अट पूर्ण झाली तर आणि तरच खरे असते:
\begin{align*}
  & \bexists{g<p}{\bexists{d<p}{\bexists{a<p}{\bexists{x<p}{\bexists{t<p}{\quad}}}}}\\
  & \qquad \fn{Sequent}(\tuple{g,d}) \land \fn{Frm}(a) \land \fn{Var}(x) \land {}\\
  & \qquad \fn{ClTerm}(t) \land {}\\
  & \qquad \fn{EndSequent}(p) =
  \tuple{
    g,
    d \concat \tuple{\Gn{\lexists} \concat x \concat a}
  } \land {}\\
  & \qquad \fn{EndSequent}((p)_1) =
  \tuple{
    g,
    d \concat \tuple{\fn{Subst}(a, t, x)}
  } \land {}\\
  & \qquad (p)_0 = 1 \land \fn{LastRule}(p) = 18.
\end{align*}
प्रस्तावित निष्पत्तीगाठीचा संकेतांक नेमक्या पुढीलपैकी एका रूपात असावा:
\begin{align*}
\fn{ProofNode}(p) \defiff {}& p=\tuple{0,\fn{EndSequent}(p)} \lor {}\\
& p=\tuple{1,(p)_1,\fn{EndSequent}(p),\fn{LastRule}(p)} \lor {}\\
& p=\tuple{2,(p)_1,(p)_2,\fn{EndSequent}(p),\fn{LastRule}(p)}.
\end{align*}
ही अट अनावश्यक अतिरिक्त घटक किंवा चुकीची गाठीची स्थानसंख्या वगळते.
मग $\fn{Correct}(p)$ ची व्याख्या अशी करतो:
\begin{multline*}
  \fn{ProofNode}(p) \land \fn{Sequent}(\fn{EndSequent}(p)) \land {}\\
  [(\fn{LastRule}(p) = 1 \land
  \fn{FollowsBy}_{\LeftR\Weakening}(p)) \lor \dots \lor {}\\
  (\fn{LastRule}(p) = 20 \land \fn{FollowsBy}_{\eq}(p)) \lor {}\\
  (p)_0 = 0 \land \fn{InitialSeq}(\fn{EndSequent}(p))]
\end{multline*}
पहिली ओळ गाठीच्या संकेतांकाचे नेमके रूप तपासते आणि $p$ ची अंतिम क्रमवर्ती खरोखर वाक्यांचीच
क्रमवर्ती आहे, याची खात्री करते. शेवटची ओळ $p$ मध्ये फक्त आरंभीची
क्रमवर्ती असण्याचा प्रसंग हाताळते.

\end{proof}

\begin{prob}
  \ref{inc:art:plk:prop:followsby} प्रमाणे पुढील गुणधर्म
  परिभाषित करा:
  \begin{enumerate}
  \item $\fn{FollowsBy}_{\Cut}(p)$,
  \item $\fn{FollowsBy}_{\LeftR\lif}(p)$,
  \item $\fn{FollowsBy}_{\eq}(p)$,
  \item $\fn{FollowsBy}_{\RightR{\lforall}}(p)$.
  \end{enumerate}
  शेवटच्या गुणधर्मासाठी, गोडेल संख्या~$p$ असलेल्या
  निष्पत्तीमधील अखेरचे निगमन आयगेन चराची अट पूर्ण करते
  का, हे आदिम पुनरावर्ती रीतीने तपासता येते, हेही दाखवावे लागेल.
  म्हणजेच $\RightR{\lforall}$ चा आयगेन चर~$a$ अंतिम
  क्रमवर्तीत येत नाही.
\end{prob}

\begin{prop}
  \label{inc:art:plk:prop:deriv}
  $p$ ही योग्य निष्पत्ती~$\pi$ ची गोडेल संख्या असेल, तर खरा
  असणारा $\fn{Deriv}(p)$ संबंध आदिम पुनरावर्ती आहे.
\end{prop}

\begin{proof}
  निष्पत्ती~$\pi$ मधील प्रत्येक निगमन एखाद्या नियमाचा योग्य
  वापर असेल, म्हणजेच तिच्या प्रत्येक उप-निष्पत्त्यांचा शेवट
  योग्य निगमनाने होत असेल, तर ती योग्य असते. म्हणून
  $\fn{Deriv}(p)$ तर आणि तरच
  \[  \bforall{i<\len{\fn{SubtreeSeq}(p)}}{\fn{Correct}((\fn{SubtreeSeq}(p))_i)}.
  \]
\end{proof}

\begin{prop}
समजा $\Gamma$ हा वाक्यांचा आदिम पुनरावर्ती संच आहे.
मग $\Prf[\Gamma](x, y)$ हा संबंध आदिम पुनरावर्ती आहे. हा संबंध
असे व्यक्त करतो: $x$ हा $\Gamma_0 \Sequent \varphi$ च्या
निष्पत्ती~$\pi$ चा संकेतांक आहे, येथे काही सांत
$\Gamma_0 \subseteq \Gamma$ आहे, आणि $y$ ही $\varphi$ ची गोडेल संख्या आहे.
\end{prop}

\begin{proof}
समजा ``$y \in \Gamma$'' हे $R_\Gamma(y)$ या आदिम पुनरावर्ती
विधेयाने दिले आहे. $\Prf[\Gamma](x, y)$ हा संबंध आदिम पुनरावर्ती आहे,
हे दाखवायचे आहे. $y$ ही वाक्य~$\varphi$ ची गोडेल संख्या असेल
आणि $x$ हा अंतिम क्रमवर्ती $\Gamma_0 \Sequent \varphi$ असलेल्या
$\Log{LK}$-निष्पत्तीचा संकेतांक असेल, तर आणि तरच हा संबंध खरा असतो.

मागील प्रस्तावानुसार, $x$ हा $\Log{LK}$ मधील योग्य
निष्पत्ती~$\pi$ चा संकेतांक असेल, तर आणि तरच खरा असणारा
$\fn{Deriv}(x)$ गुणधर्म आदिम पुनरावर्ती आहे. $x$ असा संकेतांक असेल,
तर $\fn{EndSequent}(x)$ हा $\pi$ च्या अंतिम क्रमवर्तीचा संकेतांक आहे.
म्हणून $(\fn{EndSequent}(x))_0$ हा तिच्या डाव्या बाजूचा आणि
$(\fn{EndSequent}(x))_1$ हा उजव्या बाजूचा संकेतांक आहे. त्यामुळे
``$\pi$ च्या अंतिम क्रमवर्तीची उजवी बाजू $\varphi$ आहे'' हे
$\len{(\fn{EndSequent}(x))_1} = 1 \land ((\fn{EndSequent}(x))_1)_0 =
y$ असे व्यक्त करता येते. $\pi$ च्या अंतिम क्रमवर्तीची डावी बाजू
अर्थातच सांत असते; तिच्यातील प्रत्येक वाक्य $\Gamma$ मध्ये आहे,
हे व्यक्त करणे उरते. म्हणून $\Prf[\Gamma](x, y)$ ची व्याख्या अशी करता येते:
\begin{align*}
\Prf[\Gamma](x, y) \defiff {}&
\fn{Deriv}(x) \land {} \\
& \bforall{i <
  \len{(\fn{EndSequent}(x))_0}}{R_\Gamma(((\fn{EndSequent}(x))_0)_i)} \land {}\\
& \len{(\fn{EndSequent}(x))_1} = 1 \land ((\fn{EndSequent}(x))_1)_0 = y.
\end{align*}
\end{proof}











\section{नैसर्गिक निगमनातील निष्पत्ती}\label{inc:art:pnd:sec}

\begin{explain}
निष्पत्त्यांचे अंकगणितीकरण करण्यासाठी निष्पत्त्या संख्यांनी
दर्शवावे लागतात. प्रत्येक निगमनाला एक किंवा दोन लेबले असलेल्या
सूत्रांचे वृक्ष म्हणजे निष्पत्त्या असल्यामुळे, पुनरावर्ती
निरूपण हा सहज मार्ग आहे. आपण निष्पत्ती एका क्रमित समूहाने दर्शवतो.
त्याचे घटक म्हणजे अखेरच्या निगमनाच्या आधारविधानांकडे नेणाऱ्या तात्काळ
उप-निष्पत्त्यांची संख्या, त्या उप-निष्पत्त्यांची निरूपणे,
अंतिम सूत्र, अखेरच्या निगमनाचे मुक्तता-लेबल आणि त्या निगमनाचा
प्रकार दर्शवणारी संख्या.
\end{explain}

\begin{defn}
  $\delta$ ही नैसर्गिक निगमनातील निष्पत्ती असेल, तर
  $\Gn{\delta}$ ची विगमनाने व्याख्या पुढीलप्रमाणे करतो:
  \begin{enumerate}
  \item
    $\delta$ मध्ये फक्त गृहीतक~$\varphi$ असेल, तर $\Gn{\delta}$ म्हणजे
    $\tuple{0, \Gn{\varphi}, n}$. हे गृहीतक मुक्त न केलेले असेल,
    तर~$n$ ही संख्या~$0$ असते; अन्यथा ती त्याचे संख्यात्मक लेबल असते.
  \item
    $\delta$ चा शेवट शून्य, एक, दोन किंवा तीन आधारविधानांच्या
    निगमनाने होत असेल, तर $\Gn{\delta}$ अनुक्रमे
    \begin{align*}
      & \tuple{0, \Gn{\varphi}, n, k}, \\
      & \tuple{1, \Gn{\delta_1}, \Gn{\varphi}, n, k}, \\
      & \tuple{2, \Gn{\delta_1}, \Gn{\delta_2}, \Gn{\varphi}, n, k}, \text{ किंवा}\\
      & \tuple{3, \Gn{\delta_1}, \Gn{\delta_2}, \Gn{\delta_3}, \Gn{\varphi}, n,
        k},
    \end{align*}
    असते. येथे $\delta_1$, $\delta_2$, $\delta_3$ या $\delta$ मधील
    अखेरच्या निगमनाच्या आधारविधानांवर संपणाऱ्या उप-निष्पत्त्या आहेत;
    $\varphi$ हा $\delta$ मधील त्या निगमनाचा निष्कर्ष आहे; $n$ हे त्या निगमनाचे
    मुक्तता-लेबल आहे (निगमनाने एकही गृहीतक मुक्त होत नसेल, तर~$0$);
    आणि अखेरच्या निगमनात वापरलेल्या नियमानुसार $k$ पुढील तक्त्यात
    दिले आहे.

    \begin{tabular}{lccccccc}
      \text{नियम:} & \Intro{\land} & \Elim{\land} & \Intro{\lor} & \Elim{\lor} \\
      $k$: & 1 & 2 & 3 & 4  \\[1.5ex]
      \text{नियम:} &  \Intro{\lif} & \Elim{\lif} & \Intro{\lnot} & \Elim{\lnot} \\
      $k$: & 5 & 6 & 7 & 8 \\[1.5ex]
      \text{नियम:}  &  \FalseInt & \FalseCl & \Intro{\lforall} & \Elim{\lforall} \\
      $k$: & 9 & 10 & 11 & 12 \\[1.5ex]
      \text{नियम:} & \Intro{\lexists} & \Elim{\lexists} & \Intro{\eq} & \Elim{\eq} \\
      $k$: & 13 & 14 & 15 & 16
    \end{tabular}
  \end{enumerate}
\end{defn}
\begin{ex}
  पुढील अगदी सोपी निष्पत्ती पाहा:
  \begin{prooftree}
    \AxiomC{$\Discharge{\varphi \land \psi}{1}$}
    \RightLabel{\Elim{\land}}
    \UnaryInfC{$\varphi$}
    \DischargeRule{\Intro{\lif}}{1}
    \UnaryInfC{$(\varphi \land \psi) \lif \varphi$}
  \end{prooftree}
  गृहीतकाची गोडेल संख्या $d_0 = \tuple{0,
    \Gn{(\varphi \land \psi)}, 1}$ असेल. $\Elim{\land}$ च्या निष्कर्षावर
  संपणाऱ्या निष्पत्तीची गोडेल संख्या $d_1 = \tuple{1,
     d_0, \Gn{\varphi}, 0, 2}$ असेल. येथे $1$ हे $\Elim{\land}$ ला एकच
  आधारविधान असल्यामुळे आले आहे; निष्कर्ष~$\varphi$ ची गोडेल संख्या पुढे
  आहे; $0$ हे एकही गृहीतक मुक्त होत नसल्यामुळे आले आहे; आणि $2$
  हा $\Elim{\land}$ चा संकेतांक आहे. मग संपूर्ण निष्पत्तीची
  गोडेल संख्या $\tuple{1, d_1, \Gn{((\varphi \land \psi) \lif
      \varphi)}, 1, 5}$, म्हणजेच
  \[
  \tuple{1, \tuple{1, \tuple{0, \Gn{(\varphi \land \psi)}, 1}, \Gn{\varphi}, 0, 2},
    \Gn{((\varphi \land \psi) \lif \varphi)}, 1, 5}.
  \]
\end{ex}

\begin{explain}
निष्पत्त्यांचे निरूपण ठरल्यानंतर, अशा निष्पत्त्यांच्या
गोडेल संख्यांवर आदिम पुनरावर्ती रीतीने क्रिया करता येतात आणि त्यांचे
आवश्यक गुणधर्म व संबंध व्यक्त करता येतात, हेही दाखवावे लागेल.
काही क्रिया सोप्या आहेत. उदाहरणार्थ, निष्पत्तीची गोडेल
संख्या~$d$ दिली असता, $\fn{EndFmla}(d) = (d)_{(d)_0+1}$ हे
तिच्या अंतिम सूत्राची गोडेल संख्या देते;
$\fn{DischargeLabel}(d) = (d)_{(d)_0 + 2}$ हे मुक्तता-लेबल देते;
आणि $\fn{LastRule}(d) = (d)_{(d)_0 + 3}$ हे अखेरच्या निगमनाचा
प्रकार दर्शवणारी संख्या देते. काही इतर क्रिया अधिक कठीण आहेत;
त्या कशा कराव्यात याची किमान रूपरेषा आपण पाहू. $\delta$ ही
$\Gamma$ पासून $\varphi$ ची निष्पत्ती आहे, हा संबंध
$\delta$ आणि $\varphi$ यांच्या गोडेल संख्यांचा आदिम पुनरावर्ती संबंध
आहे, हे दाखवणे आपले उद्दिष्ट आहे.
\end{explain}

\begin{prop}
  पुढील संबंध आदिम पुनरावर्ती आहेत:
  \begin{enumerate}
  \item $\varphi$ हे $\delta$ मध्ये लेबल~$n$ असलेले गृहीतक म्हणून येते.
  \item $\delta$ मधील लेबल~$n$ असलेली सर्व गृहीतके $\varphi$ या रूपाची
    असतात (म्हणजे $\delta$ मध्ये लेबल~$n$ वापरून गृहीतक~$\varphi$
    मुक्त करता येते).
  \end{enumerate}
\end{prop}

\begin{proof}
  सूत्रांच्या गोडेल संख्या आणि निष्पत्त्यांच्या गोडेल
  संख्या यांच्यातील अनुरूप संबंध आदिम पुनरावर्ती आहेत, हे दाखवायचे आहे.
  \begin{enumerate}
  \item $\fn{Assum}(x, d, n)$ आदिम पुनरावर्ती आहे, हे दाखवू.
    गोडेल संख्या~$d$ असलेल्या निष्पत्तीमधील लेबल~$n$ च्या
    गृहीतकाची गोडेल संख्या~$x$ असेल, तर हा संबंध खरा असतो.
    गोडेल संख्या~$\tuple{0, x, n}$ असलेली निष्पत्ती ही
    $d$ ची उप-निष्पत्ती असेल, तर आणि तरच असे घडते. येथे
    निष्पत्त्यांचे आपले सांकेतीकरण हे
    \ref{cmp:rec:tre:sec} मध्ये दिलेल्या वृक्षांच्या
    सांकेतीकरणाचे विशेष प्रकरण आहे, हे लक्षात घ्या. म्हणून आदिम
    पुनरावर्ती फलन $\fn{SubtreeSeq}(d)$ हे $d$ च्या सर्व
    उप-निष्पत्त्यांच्या गोडेल संख्यांची क्रमिका देते
    ($d$ च्या या क्रमिकेत पुनरुक्तीही असू शकतात). म्हणून व्याख्या अशी करता येते:
    \[    \fn{Assum}(x, d, n) \defiff \bexists{i<\len{\fn{SubtreeSeq}(d)}}{(\fn{SubtreeSeq}(d))_i =
      \tuple{0, x, n}}.
    \]
  \item $\fn{Discharge}(x, d, n)$ आदिम पुनरावर्ती आहे, हे दाखवू.
    गोडेल संख्या~$d$ असलेल्या निष्पत्तीमधील लेबल~$n$
    असलेली सर्व गृहीतके गोडेल संख्या~$x$ असलेल्या सूत्र
    चीच असतील, तर हा संबंध खरा असतो. तो
    $\bforall{y<d}{(\fn{Assum}(y, d, n) \lif y
      = x)}$ असेल, तर आणि तरच खरा असतो.
  \end{enumerate}
\end{proof}

\begin{prop}
  \label{inc:art:pnd:prop:followsby}
  गोडेल संख्या~$d$ असलेल्या निष्पत्ती~$\delta$ मधील अखेरचे
  निगमन योग्य असेल, तर आणि तरच खरा असणारा $\fn{Correct}(d)$ गुणधर्म
  आदिम पुनरावर्ती आहे.
\end{prop}

\begin{proof}
  प्रत्येक निगमन नियम~$R$ साठी $\fn{FollowsBy}_R(d)$ हा संबंध
  आदिम पुनरावर्ती आहे, हे दाखवावे लागेल. $d$ ही निष्पत्ती~$\delta$
  ची गोडेल संख्या असेल आणि $\delta$ चे अंतिम सूत्र $\delta$ च्या
  तात्काळ उप-निष्पत्त्या पासून $R$ च्या योग्य वापराने
  निष्पन्न होत असेल, तर आणि तरच $\fn{FollowsBy}_R(d)$ खरे असते.

  \Intro{\land} नियमाचे उदाहरण सोपे आहे. $\delta$ चा शेवट
  $\Intro{\land}$ च्या योग्य वापराने होत असेल, तर तिचे रूप असे असते:
  \begin{prooftree}
    \AxiomC{}
    \RightLabel{$\delta_1$}
    \DeduceC{$\varphi$}

    \AxiomC{}
    \RightLabel{$\delta_2$}
    \DeduceC{$\psi$}

    \RightLabel{\Intro\land}
    \BinaryInfC{$\varphi \land \psi$}
  \end{prooftree}
  मग $\delta$ ची गोडेल संख्या~$d$ ही
  $\tuple{2, d_1, d_2, \Gn{(\varphi \land \psi)}, 0, k}$ असते. येथे
  $\fn{EndFmla}(d_1) = \Gn{\varphi}$,
  $\fn{EndFmla}(d_2) = \Gn{\psi}$, $n=0$ आणि $k=1$.
  म्हणून $\fn{FollowsBy}_{\Intro\land}(d)$ ची व्याख्या अशी करता येते:
  \begin{multline*}
    (d)_0 = 2 \land \fn{DischargeLabel}(d) = 0 \land \fn{LastRule}(d) = 1 \land {}\\
    \fn{EndFmla}(d) = {}\\
    \Gn{(} \concat \fn{EndFmla}((d)_1) \concat \Gn{\land}
    \concat \fn{EndFmla}((d)_2) \concat \Gn{)}.
  \end{multline*}

  आणखी एक सोपे उदाहरण $\Intro\eq$ नियमाचे आहे. त्याला
  आधारविधाने नसतात, म्हणून गृहीतकांप्रमाणेच $(d)_0 = 0$ असते.
  त्याचे मुक्तता-लेबलही नसते, म्हणजे $n=0$. मात्र $\varphi$ हे
  एखाद्या बंद पद~$t$ साठी $\eq[t][t]$ या रूपाचे असले पाहिजे.
  येथे आदिम पुनरावर्ती व्याख्या अशी आहे:
  \begin{multline*}
    (d)_0 = 0 \land \fn{DischargeLabel}(d) = 0 \land {}\\
    \bexists{t<d}{(\fn{ClTerm}(t) \land
      \fn{EndFmla}(d) = {}\\
      \Gn{{\eq}(} \concat t \concat \Gn{,} \concat t \concat \Gn{)})}.
  \end{multline*}

  अधिक गुंतागुंतीचे उदाहरण म्हणून $\fn{FollowsBy}_{\Intro{\lif}}(d)$
  पाहू. $\delta$ चे अंतिम सूत्र $(\varphi \lif \psi)$ या रूपाचे,
  $\delta_1$ चे अंतिम सूत्र~$\psi$ असेल आणि $\delta$ मधील
  लेबल~$n$ असलेले प्रत्येक गृहीतक $\varphi$ या रूपाचे असेल,
  तर आणि तरच हा संबंध खरा असतो. हे आदिम पुनरावर्ती रीतीने
  पुढीलप्रमाणे व्यक्त करता येते:
  \begin{multline*}
    (d)_0 = 1 \land {}\\
    \bexists{a<d}{(\fn{Discharge}(a, (d)_1, \fn{DischargeLabel}(d)) \land {}}\\
      \fn{EndFmla}(d) = (\Gn{(} \concat a \concat \Gn{\lif}
      \concat \fn{EndFmla}((d)_1) \concat \Gn{)}))
  \end{multline*}
  (येथे $a$ ही $\varphi$ ची गोडेल संख्या समजा.)

  आणखी एका उदाहरणासाठी \Intro{\lexists} पाहू. $\delta$ मधील
  अखेरचे निगमन योग्य असेल, तर आणि तरच सूत्र~$\varphi$,
  बंद पद~$t$ आणि चल~$x$ अशी असतात की
  $\Subst{\varphi}{t}{x}$ हे निष्पत्ती~$\delta_1$ चे अंतिम
  सूत्र असते आणि $\lexists[x][\varphi]$ हा अखेरच्या निगमनाचा
  निष्कर्ष असतो. म्हणून $\fn{FollowsBy}_{\Intro{\lexists}}(d)$
  पुढील अट पूर्ण झाली तर आणि तरच खरे असते:
  \begin{multline*}
    (d)_0 = 1 \land \fn{DischargeLabel}(d) = 0 \land {} \\
    \bexists{a < d}{\bexists{x<d}{\bexists{t<d}{
          (\fn{ClTerm}(t) \land \fn{Var}(x)  \land {}}}}\\
    \fn{Subst}(a,t,x) = \fn{EndFmla}((d)_1) \land
    \fn{EndFmla}(d) = (\Gn{\lexists} \concat x \concat a)).
  \end{multline*}

नेमका गाठीचा आकार तपासण्यासाठी मागील विभागातील वैध क्रमिकासंकेतांकाची
अट वापरा. गृहीतकाला तीन घटक असतात; नियमगाठीला आधारविधानांच्या
संख्येपेक्षा चार अधिक घटक असतात:
\[\fn{NDNode}(d) \defiff \fn{Seq}(d) \land
[((d)_0=0 \land \len{d}=3) \lor ((d)_0\le3 \land \len{d}=(d)_0+4)].
\]
  मग $\fn{Correct}(d)$ ची व्याख्या अशी करतो:
  \begin{multline*}
    \fn{NDNode}(d) \land \fn{Sent}(\fn{EndFmla}(d)) \land {}\\
    [(\fn{LastRule}(d) = 1 \land
    \fn{FollowsBy}_{\Intro\land}(d)) \lor \dots \lor {}\\
    (\fn{LastRule}(d) = 16 \land \fn{FollowsBy}_{\Elim\eq}(d)) \lor {}\\
    \bexists{n<d}{\bexists{x<d}{(d = \tuple{0, x, n})}}].
  \end{multline*}
  पहिली ओळ संकेतांकाचा नेमका गाठीचा आकार आणि $d$ चे अंतिम सूत्र वाक्य आहे, याची खात्री
  करते. शेवटची ओळ $d$ मध्ये फक्त गृहीतक असण्याचा प्रसंग हाताळते.
\end{proof}

\begin{prob}
  \ref{inc:art:pnd:prop:followsby} प्रमाणे पुढील गुणधर्म
  परिभाषित करा:
  \begin{enumerate}
  \item $\fn{FollowsBy}_{\Elim{\lif}}(d)$,
  \item $\fn{FollowsBy}_{\Elim{\eq}}(d)$,
  \item $\fn{FollowsBy}_{\Elim{\lor}}(d)$,
  \item $\fn{FollowsBy}_{\Intro{\lforall}}(d)$.
  \end{enumerate}
  शेवटच्या गुणधर्मासाठी, गोडेल संख्या~$d$ असलेल्या
  निष्पत्तीमधील अखेरचे निगमन आयगेन चराची अट पूर्ण करते
  का, हे आदिम पुनरावर्ती रीतीने तपासता येते, हेही दाखवा.
  म्हणजेच $\Intro{\lforall}$ निगमनाचा आयगेन चर~$a$ हा
  $d$ च्या अंतिम सूत्रामध्ये किंवा $d$ च्या मुक्त न केलेल्या
  गृहीतकात येत नाही. यासाठी
  \ref{inc:art:pnd:prop:openassum} मधील आदिम पुनरावर्ती
  विधेय $\fn{OpenAssum}$ वापरू शकता.
\end{prob}

\begin{prop}
  \label{inc:art:pnd:prop:deriv}
  $d$ ही योग्य निष्पत्ती~$\delta$ ची गोडेल संख्या असेल,
  तर खरा असणारा $\fn{Deriv}(d)$ संबंध आदिम पुनरावर्ती आहे.
\end{prop}

\begin{proof}
  निष्पत्ती~$\delta$ मधील प्रत्येक निगमन एखाद्या नियमाचा
  योग्य वापर असेल, म्हणजेच तिच्या प्रत्येक उप-निष्पत्त्यांचा
  शेवट योग्य निगमनाने होत असेल, तर ती योग्य असते. म्हणून
  $\fn{Deriv}(d)$ तर आणि तरच
  \[  \bforall{i<\len{\fn{SubtreeSeq}(d)}}{\fn{Correct}((\fn{SubtreeSeq}(d))_i)}
  \]
\end{proof}

\begin{prop}
  \label{inc:art:pnd:prop:openassum} $z$ ही गोडेल संख्या
  मुक्त न केलेले गृहीतक~$\varphi$ ची असेल, आणि हे गृहीतक
  गोडेल संख्या~$d$ असलेल्या निष्पत्ती~$\delta$ मध्ये असेल,
  तर खरा असणारा $\fn{OpenAssum}(z, d)$ संबंध आदिम पुनरावर्ती आहे.
\end{prop}

\begin{proof}
  गृहीतकाच्या एखाद्या आस्थितीला शून्येतर लेबल~$n$ असेल आणि ती
  $\delta$ च्या अशा उप-निष्पत्तीमध्ये येत असेल जिचा शेवट
  मुक्तता-लेबल~$n$ असलेल्या नियमाने होतो, तर ती आस्थिती
  मुक्त केलेले असते. म्हणून अशा गृहीतकाची किमान एक आस्थिती
  $\delta$ मध्ये मुक्त केलेले नसेल, तर $\varphi$ हे
  $\delta$ चे मुक्त न केलेले गृहीतक असते. येथे सावध राहिले
  पाहिजे: $\delta$ मध्ये $\varphi$ च्या मुक्त केलेले आणि
  मुक्त न केलेले अशा दोन्ही आस्थिती असू शकतात.

  $\delta_0$, \dots, $\delta_k$ ही क्रमिका पाहा. येथे $\delta_0 =
  \delta$, $\delta_k$ हे गृहीतक $\Discharge{\varphi}{n}$ आहे
  (एखाद्या~$n$ साठी), आणि $\delta_{i+1}$ ही $\delta_i$ ची तात्काळ
  उप-निष्पत्ती आहे. लेबल शून्य असेल, तर हे गृहीतक व्याख्येनेच
  मुक्त न केलेले असते. अन्यथा, अशा क्रमिकेत गृहीतकाच्या आधीच्या कोणत्याही
  $\delta_i$ चा शेवट मुक्तता-लेबल~$n$ असलेल्या निगमनाने होत नसेल,
  तर $\varphi$ हे $\delta$ चे मुक्त न केलेले गृहीतक आहे.

  आदिम पुनरावर्ती फलन $\fn{SubtreeSeq}(d)$ हे $\delta$ च्या
  सर्व उप-निष्पत्त्यांच्या गोडेल संख्यांची क्रमिका देते.
  $\delta$ च्या मूळापासून गृहीतकापर्यंतचा उप-निष्पत्त्यांचा मार्ग त्यातून
  निवडता येतो. अशी उपक्रमिका ओळखणे हा आदिम पुनरावर्ती संबंध आहे:
  $\fn{Subseq}(s, s')$ तर आणि तरच
  $\bforall{i<\len{s}}{\lexists[j<\len{s'}][(s)_i = (s')_j]}$.
  येथे दिलेली अट प्रत्येक घटकाचे सदस्यत्व तपासते; ती स्वतःहून
  क्रम-संरक्षण व्यक्त करत नाही. तात्काळ उप-निष्पत्ती असण्याचा
  संबंधही आदिम पुनरावर्ती आहे: $\fn{Subderiv}(d, d')$ तर आणि तरच
  $\bexists{j<(d')_0}{d = (d')_{j+1}}$. म्हणून
  $\fn{OpenAssum}(z, d)$ ची व्याख्या अशी करता येते:
  \begin{multline*}
    \bexists{s<1+\fn{sequenceBound}(d,d)}{(\fn{Seq}(s) \land \len{s}>0 \land {}\\
      \fn{Subseq}(s, \fn{SubtreeSeq}(d))
      \land (s)_0 = d \land {}} \\
    \bexists{n<d}{((s)_{\len{s} \tsub 1} = \tuple{0, z, n} \land {}}\\
      \bforall{i<(\len{s} \tsub 1)}{(\fn{Subderiv}((s)_{i+1}, (s)_i) \land {}}\\
      (n=0 \lor \fn{DischargeLabel}((s)_i) \neq n)))).
  \end{multline*}
मार्गात प्रत्येक पुढचा गाठीचा संकेतांक आधीच्यापेक्षा लहान असतो.
म्हणून वैध मार्गात जास्तीत जास्त $d$ गाठी आणि प्रत्येक संकेतांक
जास्तीत जास्त $d$ असतो; वरची आदिम पुनरावर्ती मर्यादा पुरेशी आहे.

\end{proof}

\begin{prop}
  \label{inc:art:pnd:prop:prf-prim-rec}
  समजा $\Gamma$ हा वाक्यांचा आदिम पुनरावर्ती संच आहे.
  मग $\Prf[\Gamma](x, y)$ हा संबंध आदिम पुनरावर्ती आहे. तो असे व्यक्त
  करतो: $x$ हा $\Gamma$ मधील मुक्त न केलेले गृहीतकांपासून
  $\varphi$ च्या निष्पत्ती~$\delta$ चा संकेतांक आहे, आणि
  $y$ ही $\varphi$ ची गोडेल संख्या आहे.
\end{prop}

\begin{proof}
  समजा ``$y \in \Gamma$'' हे $R_\Gamma(y)$ या आदिम पुनरावर्ती
  विधेयाने दिले आहे. $\Prf[\Gamma](x, y)$ आदिम पुनरावर्ती आहे,
  हे दाखवायचे आहे. $y$ ही वाक्य~$\varphi$ ची गोडेल संख्या असेल,
  आणि $x$ हा अंतिम सूत्र~$\varphi$ असलेल्या नैसर्गिक निगमनातील
  निष्पत्तीचा संकेतांक असेल, तसेच तिची सर्व
  मुक्त न केलेले गृहीतके $\Gamma$ मध्ये असतील, तर आणि तरच हा
  संबंध खरा असतो.

  \ref{inc:art:pnd:prop:deriv} नुसार, $x$ ही नैसर्गिक निगमनातील योग्य
  निष्पत्ती~$\delta$ ची गोडेल संख्या असेल, तर आणि तरच खरा
  असणारा $\fn{Deriv}(x)$ गुणधर्म आदिम पुनरावर्ती आहे. त्यामुळे
  $\Prf[\Gamma](x, y)$ ची व्याख्या अशी करता येते:
  \begin{align*}
    \Prf[\Gamma](x, y) \defiff {}
    & \fn{Deriv}(x) \land \fn{EndFmla}(x) = y \land {} \\
    & \bforall{z < x}{(\fn{OpenAssum}(z, x) \lif R_\Gamma(z))}.
  \end{align*}
\end{proof}











\section{स्वयंसिद्धकीय निष्पत्ती}\label{inc:art:pax:sec}

\begin{explain}
स्वयंसिद्धकीय निष्पत्त्यांचे अंकगणितीकरण करण्यासाठी
निष्पत्त्या संख्यांनी दर्शवावे लागतात. निष्पत्त्या या
सूत्रांच्या क्रमिका असल्यामुळे, प्रत्येक निष्पत्ती
तिच्यातील सूत्रांच्या संकेतांकांच्या क्रमिकेचा संकेतांक
म्हणून दर्शवणे हा सहज मार्ग आहे.
\end{explain}

\begin{defn}
$\delta$ ही $\varphi_1$, \dots,~$\varphi_n$ या सूत्रांची
स्वयंसिद्धकीय निष्पत्ती असेल, तर $\Gn{\delta}$ म्हणजे
\[\tuple{\Gn{\varphi_1}, \dots, \Gn{\varphi_n}}.
\]
\end{defn}

\begin{ex}
  पुढील अगदी सोपी निष्पत्ती पाहा:
  \begin{derivation}
  1. & $\psi \lif (\psi \lor \varphi)$ \\
  2. & $(\psi \lif (\psi \lor \varphi)) \lif (\varphi  \lif (\psi \lif (\psi \lor \varphi)))$\\
  3. & $\varphi  \lif (\psi \lif (\psi \lor \varphi))$
  \end{derivation}
  या निष्पत्तीची गोडेल संख्या
  \begin{align*}
  \openTuple\,
    & \Gn{\psi \lif (\psi \lor \varphi)}, \\
    &\Gn{(\psi \lif (\psi \lor \varphi)) \lif (\varphi  \lif (\psi \lif (\psi \lor \varphi)))},\\
    & \Gn{\varphi  \lif (\psi \lif (\psi \lor \varphi))} \,\closeTuple.
  \end{align*}
  अशी असेल.
\end{ex}

\begin{explain}
निष्पत्त्यांचे निरूपण ठरल्यानंतर, अशा निष्पत्त्यांवर
आदिम पुनरावर्ती रीतीने क्रिया करता येतात आणि त्यांचे आवश्यक
गुणधर्म व संबंध तसेच व्यक्त करता येतात, हेही दाखवावे लागेल.
काही क्रिया सोप्या आहेत. उदाहरणार्थ, निष्पत्तीची गोडेल
संख्या~$d$ दिली असता, $(d)_{\len{d}-1}$ हे तिच्या अंतिम
सूत्राची गोडेल संख्या देते. काही इतर क्रिया अधिक कठीण
आहेत; त्या कशा कराव्यात याची किमान रूपरेषा आपण पाहू.
``$\delta$ ही $\Gamma$ पासून $\varphi$ ची निष्पत्ती आहे''
हा संबंध $\delta$ आणि $\varphi$ यांच्या गोडेल संख्यांचा आदिम
पुनरावर्ती संबंध आहे, हे दाखवणे आपले उद्दिष्ट आहे.
\end{explain}

\begin{prop}
\label{inc:art:pax:prop:followsby}
पुढील संबंध आदिम पुनरावर्ती आहेत:
\begin{enumerate}
\item $\varphi$ हे स्वयंसिद्धक आहे.
\item $\delta$ मधील $i$ वी ओळ विध्यात्मक अनुमानाने समर्थित आहे.
\item $\delta$ मधील $i$ वी ओळ \QR{} ने समर्थित आहे.
\item $\delta$ ही योग्य निष्पत्ती आहे.
\end{enumerate}
\end{prop}

\begin{proof}
सूत्रांच्या गोडेल संख्या आणि निष्पत्त्यांच्या
गोडेल संख्या यांच्यातील अनुरूप संबंध आदिम पुनरावर्ती आहेत,
हे दाखवायचे आहे.
\begin{enumerate}
\item स्वयंसिद्धक-रूपबंधांची यादी आपल्याला दिली आहे; $\varphi$ हे
  स्वयंसिद्धक असेल, तर ते या रूपबंधांपैकी एखाद्याच्या रूपाचे असते.
  यादी सांत असल्यामुळे, प्रत्येक रूपबंधासाठी $\varphi$ त्या रूपाचे
  आहे का हे आदिम पुनरावर्ती रीतीने तपासता येते, एवढे दाखवणे पुरेसे आहे.
  उदाहरणार्थ पुढील स्वयंसिद्धक-रूपबंध पाहा:
  \[  \psi \lif (\chi \lif \psi).
  \]
  $\psi$ आणि $\chi$ ही सूत्रे अशी असतील की `$($', $\psi$,
  `$\lif$', `$($', $\chi$, `$\lif$', $\psi$ आणि~`$))$' या
  चिन्हमाला जोडून $\varphi$ मिळते, तर $\varphi$ हा या रूपबंधाचा नमुना आहे.
  $\varphi$ ची गोडेल संख्या $n$ असल्यास हा गुणधर्म तपासता येतो:
  क्रमिकांची जोडणी आदिम पुनरावर्ती आहे आणि हा संबंध खरा असेल,
  तर $\psi$ व $\chi$ ही $\varphi$ ची उप-सूत्रे असल्यामुळे $\psi$ व $\chi$
  यांच्या गोडेल संख्या $\varphi$ च्या गोडेल संख्येपेक्षा लहान असतात.
  म्हणून व्याख्या अशी करता येते:
  \begin{multline*}
  \fn{IsAx}_{\psi \lif (\chi \lif \psi)}(n) \defiff \bexists{b< n}{
    \bexists{c < n}{(\fn{Frm}(b) \land \fn{Frm}(c) \land {}}}\\ n =
  \Gn{(} \concat b \concat \Gn{\lif} \concat \Gn{(} \concat c \concat
  \Gn{\lif} \concat b \concat \Gn{))}).
  \end{multline*}
  प्रत्येक स्वयंसिद्धक-रूपबंधासाठी अशी व्याख्या केली, तर
  त्यांच्या विकल्पाने $\fn{IsAx}(n)$ हा गुणधर्म परिभाषित होतो:
  ``$n$ ही स्वयंसिद्धकाची गोडेल संख्या आहे.''
\item $\delta$ मधील $i$ वी ओळ विध्यात्मक अनुमानाने समर्थित असेल,
  तर आणि तरच $j$ आणि $k < i$ अशा आधीच्या ओळी असतात की
  $j$ व्या ओळीवरील सूत्र~$\varphi$ असते,
  $k$ व्या ओळीवरील सूत्र $\varphi \lif \psi$ असते, आणि
  $i$ व्या ओळीवरील सूत्र~$\psi$ असते.
  \begin{multline*}
    \fn{MP}(d, i) \defiff \bexists{j < i}{\bexists{k < i}{}}\\
    (d)_k = \Gn{(} \concat (d)_j
      \concat \Gn{\lif} \concat (d)_i \concat \Gn{)}
  \end{multline*}
  परिबद्ध संख्यकीकरण, जोडणी आणि $=$ हे आदिम पुनरावर्ती
  असल्यामुळे हा संबंधही आदिम पुनरावर्ती आहे.
\item $\delta$ मधील एखादी ओळ \QR{} ने समर्थित असेल,
  तर ती $\psi \lif \lforall[x][\varphi(x)]$ या रूपाची असते;
  आधीची एखादी ओळ, कोणत्यातरी स्थिरांक~$c$ साठी,
  $\psi \lif \varphi(c)$ असते; आणि $c$ हा $\psi$ मध्ये येत नाही.
  पुढील अटी पूर्ण झाल्या तर आणि तरच असे घडते:
  \begin{enumerate}
  \item एखादे वाक्य~$\psi$ असते (येथे हा विशेष प्रसंग;
    सर्वसाधारण नियमात सूत्रही चालते), आणि
  \item $x$ हाच मुक्त चल असू शकणारे सूत्र~$\varphi(x)$ असते, ज्यासाठी
  \item ओळ $i$ मध्ये $\psi \lif \lforall[x][\varphi(x)]$ असते,
  \item आधीच्या एखाद्या ओळीत, $j < i$, कोणत्यातरी अचर~$c$ साठी
    $\psi \lif \Subst{\varphi}{c}{x}$ असते,
  \item आणि हा अचर $\psi$ मध्ये येत नाही.
  \end{enumerate}
  या अटी आदिम पुनरावर्ती रीतीने तपासता येतात. कारण $\psi$,
  $\varphi(x)$ आणि $x$ यांच्या गोडेल संख्या $i$ व्या ओळीवरील
  सूत्राच्या गोडेल संख्येपेक्षा लहान आहेत; तसेच $c$ ची गोडेल
  संख्या $j$ व्या ओळीवरील सूत्राच्या गोडेल संख्येपेक्षा लहान आहे:
  \begin{multline*}
    \fn{QR}_1(d, i) \defiff \bexists{j<i}{\bexists{b < (d)_i}{\bexists{x <
        (d)_i}{\bexists{a < (d)_i}{\bexists{c < (d)_j}{(}}}}}
    \\ \fn{Var}(x) \land \fn{Const}(c) \land {} \\
    (d)_i = \Gn{(} \concat
    b \concat \Gn{\lif} \concat \Gn{\lforall} \concat x \concat a
    \concat \Gn{)} \land {}\\
    (d)_j = \Gn{(} \concat b \concat
    \Gn{\lif} \concat \fn{Subst}(a,c,x) \concat \Gn{)} \land {}\\ \fn{Frm}(b)
    \land \fn{Frm}(a) \land {} \bforall{k <
      \len{b}}{(b)_k \neq (c)_0})
  \end{multline*}
  येथे $c$ आणि $x$ यांना अचर आणि चल यांच्या पदरूपातील
  गोडेल संख्या मानले आहे (म्हणजे त्यांचे चिन्हसंकेतांक नाहीत).
  $x$ हा $\varphi(x)$ मधील आदेशनासाठीचा चल आहे;
  $\Subst{\varphi(x)}{c}{x}$ हे वाक्य आहे का ही मूळ
  वाक्यांपुरती तपासणी सर्वसाधारण संख्यक-नियमासाठी आवश्यक नाही.
  सूत्ररूपता तपासून, $\psi$ मधील प्रत्येक चिन्ह $c$ पेक्षा
  वेगळे आहे ही अट ठेवून $c$ हा $\psi$ मध्ये येत नाही,
  याची खात्री करतो.

  \QR{} ची दुसरी आवृत्ती सरावासाठी ठेवतो.
\item $d$ ही योग्य निष्पत्तीची गोडेल संख्या असेल, तर आणि
  तरच ती रिकामी नसते आणि तिच्यातील प्रत्येक ओळ स्वयंसिद्धक
  किंवा विध्यात्मक अनुमान अथवा \QR{} ने समर्थित असते. म्हणून:
  \[  \fn{Deriv}(d) \defiff \fn{Seq}(d) \land \len{d}>0 \land
  \bforall{i < \len{d}}{(\fn{IsAx}((d)_i) \lor
    \fn{MP}(d,i) \lor \fn{QR}(d, i))}
  \]
\end{enumerate}
\end{proof}

\begin{prob}
\ref{inc:art:pax:prop:followsby} प्रमाणे पुढील संबंध
परिभाषित करा:
\begin{enumerate}
\item $\fn{IsAx}_{\varphi \lif (\psi \lif (\varphi \land \psi))}(n)$,
\item $\fn{IsAx}_{\lforall[x][\varphi(x)] \lif \varphi(t)}(n)$,
\item $\fn{QR}_{2}(d, i)$ (\QR{} च्या दुसऱ्या आवृत्तीसाठी).
\end{enumerate}
\end{prob}

\begin{prop}
समजा $\Gamma$ हा वाक्यांचा आदिम पुनरावर्ती संच आहे.
मग $\Prf[\Gamma](x, y)$ हा संबंध आदिम पुनरावर्ती आहे.
तो असे व्यक्त करतो: $x$ हा $\Gamma$ पासून $\varphi$ च्या
निष्पत्ती~$\delta$ चा संकेतांक आहे आणि $y$ ही
$\varphi$ ची गोडेल संख्या आहे.
\end{prop}

\begin{proof}
समजा ``$y \in \Gamma$'' हे $R_\Gamma(y)$ या आदिम पुनरावर्ती
विधेयाने दिले आहे. $\Prf[\Gamma](x, y)$ हा संबंध आदिम पुनरावर्ती
आहे, हे दाखवायचे आहे. $y$ ही वाक्य~$\varphi$ ची गोडेल संख्या
असेल आणि $x$ हा $\Gamma$ पासून $\varphi$ च्या
निष्पत्तीचा संकेतांक असेल, तर आणि तरच
$\Prf[\Gamma](x, y)$ खरा असतो.

मागील प्रस्तावानुसार, $x$ हा योग्य निष्पत्ती~$\delta$
चा संकेतांक असेल, तर आणि तरच खरा असणारा $\fn{Deriv}(x)$
गुणधर्म आदिम पुनरावर्ती आहे. मात्र त्या व्याख्येत
निष्पत्तीमधील ओळी समर्थित करण्याचा अतिरिक्त मार्ग म्हणून
$\Gamma$ संच विचारात घेतलेला नाही. ओळ \QR{} ने समर्थित आहे
का हे तपासणाऱ्या आदिम पुनरावर्ती कसोटीतही अचर~$c$ हा
$\Gamma$ मध्ये येऊ नये ही अट घेतलेली नाही. $\Gamma$ चा
थेट विचार करून व्याख्या सुधारता येईल; पण $\fn{Deriv}$ आणि
निगमन प्रमेय वापरणे सोपे आहे. काही सांत वाक्यांची यादी
$\psi_1$, \dots, $\psi_n \in \Gamma$ अशी असेल की
$\{\psi_1, \dots, \psi_n\} \Proves \varphi$, तर आणि तरच
$\Gamma \Proves \varphi$. निगमन प्रमेयानुसार, हे
$\Proves (\psi_1 \lif (\psi_2 \lif
\cdots (\psi_n \lif \varphi)\cdots))$ असेल तेव्हा घडते.
गोडेल संख्या~$z$ असलेले वाक्य या रूपाचे आहे का,
हे आदिम पुनरावर्ती रीतीने तपासता येते. म्हणून $x$ ला
गोडेल संख्या~$y$ असलेल्या वाक्याची \emph{$\Gamma$ पासून}
निष्पत्ती असण्याचा थेट संकेतांक मानण्याऐवजी,
$x$ हा वरच्या रूपातील आवर्ती सोपाधिक वाक्याच्या
$\emptyset$ पासूनच्या निष्पत्तीचा संकेतांक मानतो.
यामुळे खालील कसोटी सिद्धतायोग्यतेसाठी समतुल्य असली,
तरी संकेतांकाचा अर्थ सुरुवातीच्या विधानातील थेट संकेतांकापेक्षा बदलतो.

प्रथम, वाक्यांची क्रमिका दिल्यास तिच्यातील सर्व वाक्ये
पूर्वांग म्हणून आणि दिलेले वाक्य उत्तरांग म्हणून असलेले
सोपाधिक वाक्य आदिम पुनरावर्ती रीतीने तयार करता येते:
\begin{align*}
  \fn{hCond}(s, y, 0) & = y \\
  \fn{hCond}(s, y, n+1) & = \Gn{(} \concat (s)_{n} \concat \Gn{\lif}
  \concat \fn{hCond}(s, y, n) \concat \Gn{)}\\
  \fn{Cond}(s, y) & = \fn{hCond}(s, y, \len{s})\\
  \intertext{म्हणून $\Prf[\Gamma](x, y)$ ची व्याख्या अशी करता येते:}
  \Prf[\Gamma](x, y) & \defiff \fn{Sent}(y) \land
  \bexists{s < \fn{sequenceBound}(x,x)}{(\fn{Seq}(s) \land {}} \\
  &\qquad (x)_{\len{x}-1} = \fn{Cond}(s,y) \land {} \\
  &\qquad\bforall{i<\len{s}}{(s)_i \in \Gamma} \land {} \\
  &\qquad\fn{Deriv}(x)).
\end{align*}

येथे पुनरावर्तनाने क्रमिकेतील पूर्वांगे आतून बाहेर या क्रमाने
जोडली जातात; आधी लिहिलेल्या बाहेरून आतच्या यादीसाठी ती उलट क्रमाने
निवडता येतात. $s$ वरील परिबंध असा मिळतो: प्रत्येक $(s)_i$ ही
निष्पत्तीच्या अखेरच्या ओळीतील उप-सूत्राची गोडेल संख्या
आहे; म्हणजेच ती $(x)_{\len{x}-1}$ पेक्षा लहान आहे. $\psi
\in \Gamma$ या पूर्वांगांची संख्या, म्हणजे $s$ ची लांबी,
$x$ च्या अखेरच्या ओळीच्या लांबीपेक्षा कमी आहे.
\end{proof}



